\section{Background}

The emergence of agentic AI systems—autonomous software entities capable of reasoning, planning, delegating, and executing tasks—has introduced a fundamental governance challenge: how to maintain accountability when agents spawn sub-agents, delegate authority across organizational boundaries, and execute tasks that traverse multiple trust domains. This section establishes the theoretical and practical foundations for recursive spawning by examining agent lifecycle management, accountability architectures, delegation chain semantics, and hierarchical task decomposition paradigms.

\subsection{Agent Lifecycle Management in Multi-Agent Systems}

Multi-agent systems (MAS) have long been studied in distributed AI, but the advent of large language model (LLM)-powered agents has fundamentally altered the lifecycle landscape. Traditional software agents operated within deterministic, pre-defined boundaries; modern agentic systems exhibit non-deterministic behavior, emergent reasoning, and dynamic capability acquisition \cite{intelligent_delegation_arxiv2602}. This shift introduces lifecycle complexities that existing frameworks fail to adequately address.

An agent's lifecycle encompasses creation, initialization, execution, delegation, suspension, and termination. In single-agent systems, lifecycle management is straightforward—state transitions are bounded and observable. However, in multi-agent architectures, an agent may spawn child agents that operate concurrently, each potentially spawning further descendants. This recursive spawning pattern creates exponentially complex lifecycle graphs where parent agents must track not only their own state but also the states of all descendant agents they have spawned.

Recent research identifies the delegation chain as the critical missing layer in AI infrastructure \cite{delegated_authority_bxai}. Without explicit authority delegation protocols, agents operating in shared environments create conflicting commitments—a phenomenon demonstrated dramatically when three autonomous AI workflows in a healthcare system generated inconsistent customer communications resulting in \$1.2M in losses within three weeks \cite{decision_rights_bxai}. This failure mode underscores the necessity of lifecycle-aware delegation frameworks.

The OWL (Optimized Workforce Learning) architecture provides a hierarchical multi-agent framework that decouples strategic planning from specialized execution through a modular Planner-Coordinator-Worker structure \cite{owl_neurips2025}. This separation enables cross-domain transferability while maintaining lifecycle coherence: the Planner decomposes tasks, the Coordinator manages subtask states, and Workers execute within bounded authority. Such architectures establish a foundation for recursive spawning patterns, but they lack explicit mechanisms for immutable parent-chain provenance—the audit trail that traces every action back to its originating authority.

\subsection{Accountability and Provenance in Distributed Systems}

Accountability in distributed AI systems requires more than logging; it demands cryptographically verifiable, tamper-evident records that link every action to its authorizing entity \cite{zylos_accountability2026}. The Zylos Research tri-layer model establishes three integrated accountability requirements: identity and authentication (cryptographic binding of actions to agents), delegation and authorization (complete chain of authority from initiator to executor), and policy enforcement with non-repudiation (signed, tamper-evident records) \cite{zylos_accountability2026}. This model extends traditional audit logging by treating accountability as an identity-and-authority problem rather than merely a data storage problem.

Provenance—the documentation of data lineage and decision causality—is essential for regulatory compliance in high-stakes domains. Healthcare AI systems require comprehensive audit trails that capture data origins, transformations, and usage patterns throughout model training and deployment \cite{healthcare_ai_provenance}. The EU AI Act mandates transparency and human oversight for high-risk AI systems, requiring organizations to maintain records that enable retrospective investigation of errors or biases \cite{audit_trails_arxiv2601}. Blockchain-based approaches offer immutable anchoring of AI decisions, with cryptographic hashes recorded on distributed ledgers while heavy data (weights, logs) remains in high-throughput off-chain storage \cite{ledger_cphai}.

The COMPEL Framework introduces a Decision Event Model specifically designed for multi-agent provenance \cite{compel_framework}. This model captures domain elements including event identity, agent ID, session/workflow ID, and parent event linkage to enable reconstruction of hierarchical decision flows. Unlike traditional logging—which fails to capture delegation semantics and emergent behaviors arising from agent interactions—decision-event models explicitly represent the causal chain from inputs to outputs, answering who decided, what information was used, why the decision was made, when and in what state, how it was executed, and what alternatives were considered \cite{compel_framework}.

The PROV-AGENT framework extends the W3C PROV standard for agentic workflows, using the Model Context Protocol (MCP) to capture and relate agent interactions to broader workflow provenance \cite{prov_agent_arxiv2508}. PROV-AGENT enables fine-grained tracking of AI agent interactions—prompts, responses, decisions—within end-to-end pipelines, supporting hierarchical accountability by tracing outputs from individual agents up through the workflow.

Security considerations for multi-agent systems identify delegation chain integrity as a primary risk domain \cite{security_multiagent_arxiv2603}. Threat category ASI03—identity and privilege abuse in multi-agent delegation chains—specifically addresses the attack surface created when agents dynamically delegate authority. The analysis catalogs risks including policy-level remote code execution through tool coupling, data leakage via large-context probabilistic recall, and non-determinism as an assurance gap \cite{security_multiagent_arxiv2603}.

\subsection{Fixed vs. Mutable Delegation Chains}

A fundamental architectural decision in multi-agent systems concerns whether delegation chains are fixed at creation or mutable during execution. Fixed delegation chains establish immutable parent-child relationships at spawn time, ensuring that every action can be traced to its originating authority without ambiguity. Mutable delegation chains permit dynamic reassignment of authority, enabling adaptive reconfiguration but introducing significant accountability gaps.

Mutable delegation chains create the ``mutability problem'' identified in agentic systems discourse \cite{reddit_mutability}. A reputation system for a fixed entity is manageable; an entity that can fork, update, or be replaced introduces tracking complexities that traditional audit systems cannot accommodate. When an agent's authority is reassigned mid-execution, the provenance trail becomes fragmented—actions performed under the original delegation may be attributed to the new authority, or vice versa.

The UCAN (User Controlled Authorization Network) model addresses mutability by treating authority as delegated, scoped, and revocable rather than binary and static \cite{ixo_agent_accountable}. Actions are tied to explicit UCAN authorizations, enabling clear responsibility and visibility of misalignment. This model supports dynamic delegation while maintaining auditability by requiring that every authority transfer be itself a logged, attributable action.

Immutable parent chains represent the opposing design choice: once a delegation is established, it cannot be modified. Every descendant agent carries a permanent, verifiable link to its parent, creating an unbroken chain of custody from the root authority to the terminal action. This approach sacrifices runtime flexibility for provable accountability—any action can be traced, in constant time, to its originating agent.

The BX3 Framework adopts immutable parent chains as a core design principle. Under BX3, every spawned agent receives a permanent, cryptographically signed reference to its parent agent. This parent reference cannot be altered post-spawn, ensuring that recursive spawning creates an ever-growing, immutable audit tree rather than a mutable graph. The implications for accountability are significant: any disputed action can be traced, step-by-step, back to the root authority through the chain of parent references.

\subsection{Hierarchical Task Decomposition in AI}

Hierarchical task decomposition is the process of partitioning complex objectives into manageable subtasks organized in abstraction levels \cite{geeksforgeeks_hierarchical_planning}. This approach enables AI systems to make decisions at multiple levels of abstraction, from high-level goal formulation to primitive action execution. Hierarchical planning organizes tasks into hierarchical structures where higher-level goals decompose into sub-goals, further decomposed until reaching primitive actions \cite{geeksforgeeks_hierarchical_planning}.

Hierarchical Reinforcement Learning (HRL) extends this concept by learning multiple levels of policies corresponding to different abstraction layers \cite{geeksforgeeks_hrl}. HRL improves scalability, learning efficiency, and interpretability by decomposing tasks into subtasks and organizing policies hierarchically. Recent work demonstrates that explicit hierarchical decomposition is essential for scalable training of multi-turn LLM agents \cite{hyper_arxiv2602}.

Task decomposition benefits complex workflows by breaking them into manageable, automatable units that can be independently verified and audited \cite{task_decomposition_ai21}. This structure makes it easier to audit, refine, and scale systems that rely on large volumes of data or multi-step reasoning, especially those built using agentic AI workflows or multi-agent systems \cite{task_decomposition_ai21}. Hierarchical planning and prompt chaining decompose tasks into ordered layers or parallel modules, enabling LLM agents to plan, execute, and reflect on subtasks efficiently \cite{ai21_task_decomposition}.

The Hierarchical Task Network (HTN) planning approach solves complex problems by breaking them down into smaller, structured subtasks with explicit preconditions and effects \cite{htn_planning_gfg}. HTN planning concepts are increasingly applied to LLM-powered systems, combining structured hierarchical decomposition with the reasoning capabilities of modern language models \cite{htn_modern_agents_linkedin}. This convergence establishes HTN as a theoretical bridge between classical AI planning and contemporary agentic systems.

Recursive spawning in BX3 extends hierarchical task decomposition by making the decomposition process itself auditable and traceable. Each spawned agent represents a decomposed subtask, and the parent chain provides the decomposition tree structure. Unlike static HTN planning—where decomposition occurs at planning time—BX3 supports dynamic decomposition at runtime, with each spawning event recorded immutably in the execution trace.

\subsection{The BX3 Framework as Substrate for Immutable Parent Chains}

The BX3 Framework introduces a structured approach to agentic system governance that addresses the accountability gaps in existing multi-agent architectures. BX3 treats accountability as a first-class architectural concern, embedding provenance requirements into the agent spawning mechanism itself. Rather than retroactively constructing audit trails from logging data, BX3 makes every spawning event and delegation action a first-class, traceable occurrence in the agent lifecycle.

Central to BX3 is the POP Framework—Permissions, Obligations, and Prohibitions—that translates human policy into machine-executable governance \cite{bxai_governance}. Constitutional Charters pre-deployment-embed authority limits to eliminate governance gaps, while Evidence Packets provide auditable digital receipts documenting reasoning, inputs, and policy references for every AI action. Decision Gates implement real-time thresholds that automatically halt unauthorized AI actions, reducing reliance on manual review \cite{decision_rights_bxai}.

BX3's immutable parent chains establish a provenance substrate where every agent carries a permanent reference to its spawning authority. When an agent spawns a child, the parent reference is cryptographically signed and embedded in the child's identity record. This reference cannot be modified post-spawn, creating an ever-growing tree of attributable actions. The practical benefit is that any action—disputed, audited, or analyzed—can be traced to its root authority through the chain of parent references, without relying on mutable logging infrastructure.

The governance architecture in BX3 distinguishes between the data plane (execution) and the control plane (governance) \cite{bxai_governance}. The data plane hosts agent operations; the control plane hosts policy enforcement, provenance recording, and escalation protocols. This separation ensures that governance operations do not interfere with execution performance while maintaining comprehensive audit capability.

Recursive spawning under BX3 is not merely a scaling mechanism; it is a governance primitive. Each spawn event creates a new accountability boundary—a point at which authority is delegated, constraints are inherited, and provenance is recorded. The resulting structure is a hierarchical accountability tree where every node represents both a task decomposition and a delegation chain link. This dual nature aligns task decomposition semantics with accountability requirements, enabling both efficient execution and provable governance.

\subsection{Summary}

TheBackground establishes the necessity for immutable parent chains in recursive spawning architectures. Agent lifecycle management in multi-agent systems reveals that dynamic delegation creates accountability gaps that traditional logging cannot address. Accountability and provenance research demonstrates that cryptographic verification and decision-event models are essential for traceable AI operations. The fixed versus mutable delegation chain distinction clarifies the trade-offs between runtime flexibility and auditability. Hierarchical task decomposition provides the theoretical underpinning for how recursive spawning partitions complex tasks into manageable, attributable units. The BX3 Framework synthesizes these foundations into a coherent architecture where immutable parent chains serve as both execution infrastructure and accountability substrate.

% Bibliography
\begin{thebibliography}{99}
\small

% Agent lifecycle / delegation
 bibitem{intelligent_delegation_arxiv2602}
 Anonymous Authors.
 ``Intelligent AI Delegation.'' 
 \emph{arXiv:2602.11865}, 2026.

 bibitem{delegated_authority_bxai}
 BXAI-OS.
 ``Delegated Authority: The Missing Layer in Your AI Stack.'' 
 \url{https://bxaios.com/delegated-authority-the-missing-layer-in-your-ai-stack/}, 2026.

 bibitem{decision_rights_bxai}
 BXAI-OS.
 ``Your AI Has No Decision Rights.'' 
 \url{https://bxaios.com/decision-rights/}, 2026.

 bibitem{owl_neurips2025}
 Anonymous Authors.
 ``OWL: Optimized Workforce Learning for General Multi-Agent Assistance.'' 
 \emph{NeurIPS 2025}, 2025.

% Accountability / provenance
 bibitem{zylos_accountability2026}
 Zylos Research.
 ``AI Agent Accountability: Audit Trails, Attribution, and Non-Repudiation in Multi-Agent Systems.'' 
 \url{https://zylos.ai/research/2026-03-22-ai-agent-accountability-audit-trails-attribution-multi-agent-systems}, 2026.

 bibitem{healthcare_ai_provenance}
 Healthcare Technology.
 ``AI and LLM Data Provenance and Audit Trails for Healthcare Technology.'' 
 \url{https://www.onhealthcare.tech/p/ai-and-llm-data-provenance-and-audit}, 2024.

 bibitem{audit_trails_arxiv2601}
 Anonymous Authors.
 ``Audit Trails for Accountability in Large Language Models.'' 
 \emph{arXiv:2601.20727v1}, 2026.

 bibitem{ledger_cphai}
 Copenhagen AI.
 ``Ledger: Economic Layer.'' 
 \url{https://cph.ai/ledger}, 2026.

 bibitem{compel_framework}
 COMPEL Framework.
 ``Audit Trails and Decision Provenance in Multi-Agent Systems.'' 
 \url{https://www.compelframework.org/articles/audit-trails-and-decision-provenance-in-multi-agent-systems}, 2026.

 bibitem{prov_agent_arxiv2508}
 Anonymous Authors.
 ``PROV-AGENT: Unified Provenance for Tracking AI Agent Interactions in Agentic Workflows.'' 
 \emph{arXiv:2508.02866v1}, 2025.

 bibitem{security_multiagent_arxiv2603}
 Anonymous Authors.
 ``Security Considerations for Multi-agent Systems.'' 
 \emph{arXiv:2603.09002}, 2026.

% Mutable delegation
 bibitem{reddit_mutability}
 Anonymous Author.
 ``AI agents are hiring other AI agents. Nobody asked who's verifying...''
 \url{https://www.reddit.com/r/ArtificialInteligence/comments/1rftnk5/}, 2025.

% Hierarchical decomposition
 bibitem{geeksforgeeks_hierarchical_planning}
 GeeksforGeeks.
 ``Hierarchical Planning in AI.'' 
 \url{https://www.geeksforgeeks.org/artificial-intelligence/hierarchical-planning-in-ai/}, 2024.

 bibitem{geeksforgeeks_hrl}
 GeeksforGeeks.
 ``Hierarchical Reinforcement Learning (HRL) in AI.'' 
 \url{https://www.geeksforgeeks.org/artificial-intelligence/hierarchical-reinforcement-learning-hrl-in-ai/}, 2024.

 bibitem{hyper_arxiv2602}
 Anonymous Authors.
 ``HiPER: Hierarchical Reinforcement Learning with Explicit Credit Assignment.'' 
 \emph{arXiv:2602.16165}, 2026.

 bibitem{task_decomposition_ai21}
 AI21 Labs.
 ``What is Task Decomposition?'' 
 \url{https://www.ai21.com/glossary/foundational-llm/task-decomposition/}, 2024.

 bibitem{ai21_task_decomposition}
 AI21 Labs.
 ``Task Decomposition.'' 
 \url{https://www.ai21.com/glossary/foundational-llm/task-decomposition/}, 2024.

 bibitem{htn_planning_gfg}
 GeeksforGeeks.
 ``Hierarchical Task Network (HTN) Planning in AI.'' 
 \url{https://www.geeksforgeeks.org/artificial-intelligence/hierarchical-task-network-htn-planning-in-ai/}, 2024.

 bibitem{htn_modern_agents_linkedin}
 Anonymous Author.
 ``Day 9 - Hierarchical Task Networks in Modern Agent Systems.'' 
 \url{https://www.linkedin.com/pulse/day-9-hierarchical-task-networks-modern-agent-systems-marques-ylnce}, 2025.

% BX3 / accountable autonomy
 bibitem{bxai_governance}
 BXAI-OS.
 ``AI Governance.'' 
 \url{https://bxaios.com/ai-governance/}, 2026.

 bibitem{ixo_agent_accountable}
 IXO.
 ``Your Agent Did What!?'' 
 \url{https://impacts.ixo.world/your-agent-did-what/}, 2026.

 bibitem{emergentmind_htd}
 Emergent Mind.
 ``Hierarchical Task Decomposition.'' 
 \url{https://www.emergentmind.com/topics/hierarchical-task-decomposition}, 2024.

\end{thebibliography}